\documentclass[11pt,a4paper]{article}
\usepackage[a4paper,margin=2.5cm]{geometry}
\usepackage[T1]{fontenc}\usepackage[utf8]{inputenc}\usepackage[scaled=0.95]{helvet}\renewcommand{\familydefault}{\sfdefault}
\usepackage{textcomp}\usepackage[colorlinks,urlcolor=blue]{hyperref}\usepackage{xcolor}\definecolor{accent}{HTML}{C0571A}\definecolor{blue}{HTML}{2F6E9E}
\setlength{\parindent}{0pt}\setlength{\parskip}{8pt}\linespread{1.15}\pagestyle{empty}
\begin{document}
{\fontsize{24}{28}\selectfont\color{blue!60!black} Poly Citra CoolBrick}\par
{\color{blue}\rule{\textwidth}{0.8pt}}\par
Rian Sikka and Daiwik Abrol \quad|\quad Grade 9 (IGCSE)\\ Scottish High International School
\vspace{10pt}

{\large\bfseries\color{accent} Abstract}\par
India has about ten lakh government schools, and most have no air-conditioning. Research on ten million students found that every extra degree Fahrenheit of school-year heat reduced learning by about one percent in schools without cooling. At the same time, Styrofoam packaging fills landfills, orange peels from juice shops are dumped or burned, and clay bricks consume fertile topsoil and coal-fired kilns. Our project, Poly Citra CoolBrick, turns these two wastes into a brick designed to keep heat out of buildings.

We extracted D-limonene from discarded orange peels by steam distillation. This natural solvent dissolves the polystyrene in Styrofoam, releasing its trapped air and leaving a gel that we pressed into a mould and cured into a solid brick without firing. Our prototype floated, kept its shape when heated and supported a load. For the CoolBrick wall design, we added a white calcium carbonate coat to reflect sunlight and moulded air pockets to slow heat transfer.

A computer simulation of a sun-facing wall on a 41 \textdegree C Delhi summer day showed the CoolBrick surface peaking at 39 \textdegree C, compared with 52 \textdegree C for clay, and letting about eighty percent less heat into a building. We plan a four-hour sunlight test comparing CoolBrick and clay samples, along with laboratory tests of strength, water absorption and fire performance against Indian standards.

At an estimated twenty-four rupees per brick in batch production, the CoolBrick could cool classrooms without electricity, supporting SDGs 4, 11, 12 and 13.
\textbf{Simulation video:} \url{https://youtu.be/kSh3PBh7VoA}
\end{document}
